\documentclass[11pt,a4paper]{article}
\input{../template/utm-lab}

\newcommand{\labnumber}{7}
\newcommand{\labtitle}{Three Small Programs: Processes, Redirection, a Pipe}
\newcommand{\studentname}{Daniel Nicorici}
\newcommand{\studentgroup}{FAF-243}

\begin{document}

\utmtitlepage
\tableofcontents
\newpage

\plainsection{Purpose of the work}

The lecture asked three questions: how a process is born, where \code{printf} writes, and how
\code{ls | wc -l} works. The purpose of this lab is to run the answers: three short C programs
that do by hand what the shell does for a command, for a redirection and for a pipe. Each
program is compiled, run, changed in one line and run again to see what happens.

The work was done on my Arch Linux laptop with gcc 16.2.1, in an interactive bash in the folder
\code{\textasciitilde/os-lab7}.

\begin{lstlisting}
$ mkdir -p ~/os-lab7 && cd ~/os-lab7
$ gcc --version | head -1
gcc (GCC) 16.2.1 20260810
\end{lstlisting}

\section{Part 1. fork + exec + wait}

\code{run.c} performs the three steps of the shell for the command \code{ls -l}: the process
copies itself with \code{fork}, the copy becomes \code{ls} with \code{execlp}, and the original
waits for it with \code{waitpid}.

\begin{lstlisting}[style=ccode,caption={run.c}]
#include <stdio.h>
#include <stdlib.h>
#include <unistd.h>
#include <sys/wait.h>

/* the three steps of the shell for the command:   ls -l   */
int main(void)
{
    printf("parent: my PID is %d\n", getpid());
    fflush(stdout);

    pid_t pid = fork();                          /* step 1: copy myself */
    if (pid < 0) { perror("fork"); return 1; }

    if (pid == 0) {                              /* only the child runs this */
        printf("child:  my PID is %d, my parent is %d\n", getpid(), getppid());
        fflush(stdout);
        execlp("ls", "ls", "-l", (char *)NULL);  /* step 2: become ls */
        perror("exec");                          /* runs only if exec failed */
        exit(127);
    }

    int status;
    waitpid(pid, &status, 0);                    /* step 3: wait for the child */
    printf("parent: child %d ended with status %d\n", pid, WEXITSTATUS(status));
    return 0;
}
\end{lstlisting}

\begin{lstlisting}[caption={run: the child becomes ls -l}]
$ gcc -Wall -Wextra -o run run.c && ./run
parent: my PID is 98814
child:  my PID is 98815, my parent is 98814
total 36
-rw-r--r-- 1 Nico Nico  1013 Oct 22 20:28 pipeline.c
-rw-r--r-- 1 Nico Nico   703 Oct 22 20:28 redirect.c
-rwxr-xr-x 1 Nico Nico 16472 Oct 22 20:29 run
-rw-r--r-- 1 Nico Nico   911 Oct 22 20:28 run.c
-rw-r--r-- 1 Nico Nico   979 Oct 22 20:28 zombie.c
parent: child 98815 ended with status 0
\end{lstlisting}

Then the \code{execlp} line was replaced, first with
\code{execlp("date", "date", (char *)NULL);} and then with
\code{execlp("nosuchprog", "nosuchprog", (char *)NULL);}:

\begin{lstlisting}[caption={The child becomes date}]
$ gcc -Wall -Wextra -o run run.c && ./run
parent: my PID is 99360
child:  my PID is 99361, my parent is 99360
Thu Oct 22 20:29:21 EEST 2026
parent: child 99361 ended with status 0
\end{lstlisting}

\begin{lstlisting}[caption={A program that does not exist}]
$ gcc -Wall -Wextra -o run run.c && ./run
parent: my PID is 99460
child:  my PID is 99461, my parent is 99460
exec: No such file or directory
parent: child 99461 ended with status 127
\end{lstlisting}

\textbf{Observations:}
\begin{enumerate}
  \item \textbf{Why \code{perror("exec")} appears only with \code{nosuchprog}.} A successful
  \code{execlp} never returns: it replaces the whole program of the child with \code{ls} or
  \code{date}, so the lines after it no longer exist in that process. \code{execlp} returns only
  when it fails, and that happened only with \code{nosuchprog}, which is not found in any PATH
  directory; then \code{perror} printed the reason, ``No such file or directory''.
  \item \textbf{The child's parent.} The child prints \code{my parent is 98814}, which is the
  PID that the parent printed for itself on the first line: \code{getppid()} returns the PID of
  the process that called \code{fork}.
  \item \textbf{The status.} With \code{ls} (and \code{date}) the parent sees status 0, because
  the program ended successfully. With \code{nosuchprog} it sees 127: the exec failed, so the
  child ran its own \code{exit(127)}, and \code{WEXITSTATUS} extracted that number.
\end{enumerate}

\section{Part 2. A zombie, on purpose}

\code{zombie.c} is \code{run.c} (the \code{ls} version) with one new line before
\code{waitpid}, so the parent sleeps for 20 seconds before it collects its child:

\begin{lstlisting}[style=ccode,firstnumber=23,caption={The change in zombie.c}]
    int status;
    sleep(20);                                   /* the new line */
    waitpid(pid, &status, 0);                    /* step 3: wait for the child */
\end{lstlisting}

\begin{lstlisting}[caption={The process table while the parent sleeps}]
$ gcc -Wall -Wextra -o zombie zombie.c
$ ./zombie > /dev/null &
[1] 99669
$ sleep 2; ps -o pid,ppid,stat,cmd
    PID    PPID STAT CMD
  98636   98633 SNs  bash -i
  99669   98636 SN   ./zombie
  99670   99669 ZN   [ls] <defunct>
  99918   98636 RN+  ps -o pid,ppid,stat,cmd
\end{lstlisting}

\begin{lstlisting}[caption={The process table after 20 seconds}]
$ ps -o pid,ppid,stat,cmd
    PID    PPID STAT CMD
  98636   98633 SNs  bash -i
 100897   98636 RN+  ps -o pid,ppid,stat,cmd
[1]+  Done                       ./zombie > /dev/null
\end{lstlisting}

\textbf{Observations:}
\begin{enumerate}
  \item \textbf{The STAT letter of \code{ls}.} \code{ls} has \code{Z}: it is a zombie, shown as
  \code{[ls] <defunct>}. The child has already finished, but its parent (PID 99669,
  \code{./zombie}) has not collected its exit status yet, so the kernel keeps a small entry with
  the PID and the status in the process table. The \code{N} next to the letters only means that
  the processes of my session run with a lowered priority (a positive nice value).
  \item \textbf{After 20 seconds.} The \code{[ls] <defunct>} line is gone, and so is
  \code{./zombie}, which bash reports as \code{Done}. The line that removed the zombie is
  \code{waitpid(pid, \&status, 0);}: when the parent finally waited, it received the exit status
  and the kernel freed the zombie's entry.
\end{enumerate}

\section{Part 3. Redirection: ls -l > out.txt}

\code{redirect.c} does what the shell does for \code{ls -l > out.txt}: in the child it closes
descriptor 1, opens the file, which receives the smallest free number, and then runs \code{ls}.

\begin{lstlisting}[style=ccode,caption={redirect.c}]
#include <stdio.h>
#include <stdlib.h>
#include <fcntl.h>
#include <unistd.h>
#include <sys/wait.h>

/* does the same as the shell command:   ls -l > out.txt   */
int main(void)
{
    pid_t pid = fork();
    if (pid == 0) {
        close(1);                                /* descriptor 1 is now free */
        int fd = open("out.txt", O_WRONLY | O_CREAT | O_TRUNC, 0644);
        if (fd < 0) { perror("open"); exit(1); }
        fprintf(stderr, "open gave descriptor %d\n", fd);
        execlp("ls", "ls", "-l", (char *)NULL);  /* ls writes to 1, as always */
        perror("exec");
        exit(127);
    }
    waitpid(pid, NULL, 0);
    printf("done: the listing is in out.txt\n");
    return 0;
}
\end{lstlisting}

\begin{lstlisting}[caption={redirect and the content of out.txt}]
$ gcc -Wall -Wextra -o redirect redirect.c && ./redirect
open gave descriptor 1
done: the listing is in out.txt
$ head -3 out.txt
total 76
-rw-r--r-- 1 Nico Nico     0 Oct 22 20:29 out.txt
-rw-r--r-- 1 Nico Nico  1013 Oct 22 20:28 pipeline.c
\end{lstlisting}

\textbf{Observations:}
\begin{enumerate}
  \item \textbf{The number \code{open} gave.} It gave descriptor 1. \code{open} always returns
  the smallest descriptor that is not in use: 0 is standard input, 1 had just been closed, and 2
  is standard error, so 1 was the first free number.
  \item \textbf{Why the listing is not on the screen.} \code{ls} was not changed: it writes to
  descriptor 1, as every program does. But in this child, descriptor 1 is now \code{out.txt},
  and \code{exec} replaces the program while keeping the open descriptors, so \code{ls} writes
  into the file without knowing it. The message about descriptor 1 still appeared because it was
  written to descriptor 2, which was not touched.
  \item \textbf{Redirecting the errors.} For \code{2> err.txt}, I would close descriptor 2
  before \code{open}; the file would then become descriptor 2, standard error.
\end{enumerate}

\section{Part 4. A pipe: ls /etc | wc -l}

\code{pipeline.c} builds the pipeline by hand with one pipe and two children. In \code{ls},
descriptor 1 becomes the write end \code{p[1]}; in \code{wc}, descriptor 0 becomes the read end
\code{p[0]}. \code{dup} copies a descriptor into the smallest free number, just like
\code{open}.

\begin{lstlisting}[style=ccode,caption={pipeline.c}]
#include <stdio.h>
#include <stdlib.h>
#include <unistd.h>
#include <sys/wait.h>

/* does the same as the shell command:   ls /etc | wc -l   */
int main(void)
{
    int p[2];
    if (pipe(p) < 0) { perror("pipe"); return 1; }

    if (fork() == 0) {                           /* ls: its 1 becomes p[1] */
        close(1);
        dup(p[1]);                               /* smallest free number: 1 */
        close(p[0]);
        close(p[1]);
        execlp("ls", "ls", "/etc", (char *)NULL);
        perror("exec ls");
        exit(127);
    }

    if (fork() == 0) {                           /* wc: its 0 becomes p[0] */
        close(0);
        dup(p[0]);                               /* smallest free number: 0 */
        close(p[0]);
        close(p[1]);
        execlp("wc", "wc", "-l", (char *)NULL);
        perror("exec wc");
        exit(127);
    }

    close(p[0]);                                 /* the parent uses neither end */
    close(p[1]);
    wait(NULL);
    wait(NULL);
    return 0;
}
\end{lstlisting}

\begin{lstlisting}[caption={pipeline compared with the real shell}]
$ gcc -Wall -Wextra -o pipeline pipeline.c && ./pipeline
181
$ ls /etc | wc -l
181
\end{lstlisting}

Then the parent's \code{close(p[1]);} was turned into a comment:

\begin{lstlisting}[style=ccode,firstnumber=32,caption={The change that breaks the pipeline}]
    close(p[0]);                                 /* the parent uses neither end */
    /* close(p[1]); */
\end{lstlisting}

\begin{lstlisting}[caption={The broken pipeline never ends}]
$ gcc -Wall -Wextra -o pipeline pipeline.c && ./pipeline
^C
$ ./pipeline &
[1] 101527
$ sleep 1; ps -o pid,ppid,stat,cmd
    PID    PPID STAT CMD
  98636   98633 SNs  bash -i
 101527   98636 SN   ./pipeline
 101529  101527 SN   wc -l
 101603   98636 RN+  ps -o pid,ppid,stat,cmd
$ kill %1
[1]+  Terminated                 ./pipeline
\end{lstlisting}

The program printed nothing and had to be stopped with Ctrl-C. Running it once more in the
background showed which processes were still alive; after that, \code{pipeline.c} was restored
to the original version.

\textbf{Observations:}
\begin{enumerate}
  \item \textbf{The same number.} Yes: \code{./pipeline} and \code{ls /etc | wc -l} both
  printed 181.
  \item \textbf{Which child never ends.} \code{wc}. In the process table there is still
  \code{wc -l} (PID 101529, child of \code{./pipeline}), while \code{ls} has already finished
  and been collected by the parent's first \code{wait}.
  \item \textbf{What \code{wc} is waiting for.} \code{wc} keeps reading until end of file, and
  a pipe reports end of file only when every write end is closed, but the parent still holds
  \code{p[1]} open, so \code{wc} waits forever and the parent's second \code{wait} waits for
  \code{wc}.
\end{enumerate}

The optional three-stage \code{pipeline3.c} was not done.

\plainsection{Conclusions}

The three programs showed that the shell's everyday features are built from a few system calls.
Running a command is \code{fork}, \code{exec} and \code{wait}: the child is a copy of the parent
until \code{exec} replaces its program completely, which is why the code after a successful
\code{exec} never runs and why the exit status is the only thing the parent learns. A child that
has ended but has not been waited for stays in the process table as a zombie until the parent
calls \code{wait}.

Redirection and pipes rely on one rule: a new descriptor always gets the smallest free number.
Closing descriptor 1 and opening a file makes the file standard output, and \code{exec} keeps it,
so an unchanged \code{ls} writes into the file. A pipe connects the standard output of one child
to the standard input of another with \code{dup}, and it works only if every unused end is
closed: a single forgotten write end in the parent was enough to make \code{wc} wait forever for
an end of file that never came.

\end{document}
